\section*{Computational details}
\addcontentsline{toc}{section}{Computational details}

All results in this paper were computed with \pkg{cvxcla}~2.7.1 \citep{cvxcla} and
\proglang{Python}~3.14.6 on an Apple M4 Pro (14 cores, 48~GB of memory) running
macOS~15.6.1. \pkg{NumPy}~2.4.6 and \pkg{SciPy}~1.17.1 were linked against Apple's
Accelerate BLAS, and the first turning point under a general equality system is found
with the \pkg{HiGHS} solver through \code{scipy.optimize.linprog}. The comparisons use
\pkg{OSQP}~1.1.3, \pkg{Clarabel}~0.11.1, the \pkg{qpOASES} bundled with
\pkg{CasADi}~3.8.1, and \proglang{R}~4.6.1 with the \proglang{R} package
\pkg{CLA}~0.96-3. The estimators of Section~\ref{sec:estimators} use
\pkg{scikit-learn}~1.9.0, the figures \pkg{matplotlib}~3.11.0, and the return data are
read with \pkg{pandas}~3.0.3 and \pkg{pyarrow}~24.0.0.

\paragraph{Software availability.}
\pkg{cvxcla} is released under the MIT license. Release~2.7.1 is
archived on Zenodo \citep{cvxcla} under \doi{10.5281/zenodo.22209208}; the source is
developed at \url{https://github.com/cvxgrp/cvxcla}, and every source link in this paper
points to the release tag \code{v2.7.1}. The package declares only lower bounds on its
dependencies, and its continuous integration also tests them at those bounds; the
experiment script pins the exact versions above. The test suite, some $600$ unit and
property-based tests (the latter with \pkg{Hypothesis}), enforces full line coverage of
the library.

\paragraph{Reproducing the results.}
Every figure, table, and number is produced by one script, \code{make\_figures.py},
which the replication materials of this article hold together with the returns
snapshot. It is the experiment script of release~2.7.1
(\srcfile{experiments/make_figures.py}{make\_figures.py} in the \pkg{cvxcla} repository), less one external
timing baseline that this paper does not report and with one step added, the worked
example of Section~\ref{sec:worked}; the other steps named in the text, such as
\target{scaling}, are the same in both. The script declares the versions above as
inline (PEP~723) metadata, so with \pkg{uv} installed
\begin{CodeChunk}
\begin{CodeInput}
uv run make_figures.py all
\end{CodeInput}
\end{CodeChunk}
installs them into an isolated environment and rebuilds everything. A target name in
place of \code{all} rebuilds one artifact, and \code{all -{}-quick} skips the two scaling
sweeps (Sections~\ref{sec:scaling} and~\ref{sec:rankscaling}), which take about $45$ minutes
on the machine above. Random inputs are drawn from fixed seeds and the S\&P 500 returns
are a committed snapshot (Section~\ref{sec:empirical}), so the run needs no network access
once the packages are installed. Timings and memory were measured on this machine only and will differ elsewhere; the
turning-point counts and accuracy figures should agree up to round-off. On \proglang{Linux}
(Debian, \proglang{Python}~3.12, \pkg{NumPy}'s bundled OpenBLAS, in a container on the same
machine) the six correctness checks of Section~\ref{sec:exactness} and
Appendix~\ref{sec:details} find the same turning points; the worst KKT residual is
$7.2\times10^{-13}$ against $1.7\times10^{-13}$, and only round-off-level counts move,
such as which zero-length segments of the $130/30$ trace occur and how often the
projection of Appendix~\ref{sec:degeneracy} fires.
